% !TeX root = 00-main.tex
\documentclass[00-main.tex]{subfiles}

\begin{document}

\section{Conclusion}

This report has reviewed how fractional calculus extends classical epidemic models by introducing nonlocal temporal dependence.
Through representative compartmental models and a detailed comparison of integer- and fractional-order SEIR systems, it has shown that fractional derivatives can provide useful flexibility and may admit a stochastic interpretation.
Their use, however, requires more than improved curve fitting.
A convincing fractional epidemic model should be dimensionally consistent, connect its memory kernel to a plausible mechanism, and demonstrate that its additional parameters can be identified reliably from data.
Fractional models are therefore best regarded as carefully justified alternatives to classical models rather than automatic replacements for them.

\ifSubfilesClassLoaded{\printbibliography}{}

\end{document}
